\chapter{Intravenous Pyelogram (IVP)}
\section*{What Is This Test?} IVP is an X-ray study that uses injected contrast to display the kidneys, ureters, and bladder.\section*{Alternative Names} Intravenous urography; excretory urogram; IVU.\section*{Principle} The kidneys filter iodinated contrast, which outlines the urinary tract on timed radiographs.\section*{Why This Test Is Ordered} It may evaluate blood in urine, stones, obstruction, structural abnormalities, or urinary-tract symptoms when other imaging is unsuitable.\section*{Primary vs Secondary Test} It is a secondary imaging test; CT urography, ultrasound, or MRI may be preferred depending on the question.\section*{How This Test Is Performed} Contrast is injected into a vein and images are obtained as it passes through the urinary system.\section*{What Preparation Is Needed} Fasting or bowel preparation may be requested. Report kidney disease, diabetes, pregnancy possibility, and prior contrast reactions.\section*{Understanding Results} Findings may include obstruction, stones, delayed excretion, or structural abnormality; small lesions can be missed.\section*{Follow-up Tests} Urinalysis, CT, ultrasound, cystoscopy, or urology consultation may follow.\section*{References} \begin{enumerate}\item MedlinePlus. Intravenous Pyelogram (IVP).\item American College of Radiology. Urography imaging guidance.\end{enumerate}\clearpage\section*{Understanding Results and Follow-up}
The laboratory report should identify the specimen, method, units, reference interval or decision limit, and whether the result is positive, negative, abnormal, or inconclusive. Reference intervals vary with age, sex, pregnancy, specimen type, assay, and laboratory; a result outside a range is not automatically a diagnosis, and a result within a range does not always exclude disease. Discuss unexpected results with the ordering clinician, who can decide whether repeat or confirmatory testing, treatment, or referral is appropriate. Do not change medicines or delay urgent care based on an isolated result.
\section*{Regulatory Status and Authorization Date}This chapter describes a test category, not a specific commercial product. FDA, CDSCO, EU IVDR, and MHRA approval or registration dates therefore cannot be assigned without the assay manufacturer, product name, instrument, and intended use. For laboratory-developed tests, an individual agency approval date may not exist. See the Preface for the interpretation of regulatory status.
\clearpage
